\section{Experiments}

\subsection{Setup}

We evaluate on EvidenceBench, a synthetic benchmark containing 12,000 scientific
claims paired with abstracts from four research domains. The training, development, and
test splits contain 8,000, 2,000, and 2,000 claims, respectively. We report accuracy,
macro F1, and expected calibration error (ECE). All models use the same retrieved
evidence and are trained with three random seeds.

\subsection{Main Results}

Table~\ref{tab:main-results} compares CEA with a claim-only classifier and a verifier
that mean-pools all evidence representations. CEA obtains the strongest accuracy and
macro F1 while reducing calibration error. The improvement is largest on the subset in
which the top two evidence sources disagree.

\begin{table}[ht]
  \centering
  \caption{Test performance on EvidenceBench. Higher accuracy and macro F1 are better;
  lower ECE is better.}
  \label{tab:main-results}
  \begin{tabular}{lccc}
    \toprule
    Model & Accuracy (\%) & Macro F1 & ECE \\
    \midrule
    Claim only & 68.1 & 66.7 & 0.142 \\
    Mean pooling & 80.3 & 79.5 & 0.091 \\
    CEA & \textbf{84.6} & \textbf{83.9} & \textbf{0.047} \\
    \bottomrule
  \end{tabular}
\end{table}

\subsection{Ablation Study}

We remove one component at a time and retrain the model under the same protocol.
Removing source calibration decreases macro F1 by 3.8 points, indicating that
reliability weighting is essential when evidence sources disagree.
Removing the relevance scorer produces a smaller but consistent decline of 1.6 points.
These results suggest that the two scores capture complementary properties of the
retrieved evidence.
